% ==============================================================================
% methodology.tex
% ==============================================================================

\section{Review Methodology}
\label{sec:methodology}

This survey follows a systematic review methodology adapted from the PRISMA (Preferred Reporting Items for Systematic Reviews and Meta-Analyses) guidelines.

\subsection{Search Strategy}

We conducted literature searches across the following databases and repositories:
\begin{itemize}
  \item IEEE Xplore
  \item PubMed / MEDLINE
  \item Google Scholar
  \item arXiv (cs.CV, eess.IV)
  \item Semantic Scholar
\end{itemize}

Search queries combined terms related to X-ray-to-CT reconstruction, few-view reconstruction, single-view 3D reconstruction, and specific method families (GAN, transformer, NeRF, diffusion). The search period covered publications from January 2018 to December 2025.

\subsection{Inclusion and Exclusion Criteria}

Papers were included if they:
\begin{enumerate}
  \item Proposed or evaluated a learning-based method for reconstructing CT volumes from $\leq$10 X-ray projections.
  \item Were published in peer-reviewed venues or deposited on recognised preprint servers.
  \item Provided quantitative evaluation against established metrics (e.g., PSNR, SSIM, FID).
\end{enumerate}

Papers were excluded if they focused solely on traditional (non-learning) reconstruction, addressed sparse-view CT with $>$10 projections without relevance to the few-view setting, or were non-English publications.

\subsection{Paper Categorisation}

Selected papers were categorised into three tiers:
\begin{description}
  \item[Core papers] directly address few-view X-ray-to-CT reconstruction with a novel method or significant empirical contribution.
  \item[Supporting papers] provide relevant techniques, architectures, or insights that are applied or adapted in core methods.
  \item[Background papers] offer foundational context (e.g., classical reconstruction, general deep learning surveys).
\end{description}

% TODO: Add PRISMA flow diagram figure.
